\documentclass{article}
\usepackage{makeidx}
\makeindex

\title{pcons \LaTeX{} Example}
\author{pcons}
\date{2026}

\begin{document}
\maketitle
\tableofcontents

\section{Introduction}
\label{sec:intro}

This is a \LaTeX{} document built with pcons using the latexmk toolchain.
It exercises multi-pass compilation, bibliography processing, index generation,
cross-references, and included source files.

The build system section is in Section~\ref{sec:buildsys}.
See also the results in Section~\ref{sec:results} on page~\pageref{sec:results}.

\section{Build Systems}
\label{sec:buildsys}

\index{build system}
As shown by Knuth~\cite{knuth1984}, \TeX{} is a powerful typesetting system.
Modern build systems\index{build system!modern} like pcons~\cite{pcons2026}
generate Ninja\index{Ninja} build files for fast incremental rebuilds.

Unlike Make\index{Make}, Ninja\index{Ninja} is designed to be generated by a
higher-level tool rather than written by hand.

% This file is \input'd to test dependency tracking of included sources.
\section{Methods}
\label{sec:methods}

This section is loaded via \verb|\section{Methods}
\label{sec:methods}

This section is loaded via \verb|\section{Methods}
\label{sec:methods}

This section is loaded via \verb|\input{methods}|, which tests that
the depfile tracks included source files.  Modifying this file should
trigger a rebuild via Ninja's dependency tracking.

\index{dependency tracking}
The \texttt{-recorder} flag causes \LaTeX{} to write a \texttt{.fls} file
listing every file it reads.  The pcons helper converts this to a Ninja
depfile so that changes to any included \texttt{.tex}, \texttt{.sty},
or \texttt{.cls} file trigger a rebuild.
|, which tests that
the depfile tracks included source files.  Modifying this file should
trigger a rebuild via Ninja's dependency tracking.

\index{dependency tracking}
The \texttt{-recorder} flag causes \LaTeX{} to write a \texttt{.fls} file
listing every file it reads.  The pcons helper converts this to a Ninja
depfile so that changes to any included \texttt{.tex}, \texttt{.sty},
or \texttt{.cls} file trigger a rebuild.
|, which tests that
the depfile tracks included source files.  Modifying this file should
trigger a rebuild via Ninja's dependency tracking.

\index{dependency tracking}
The \texttt{-recorder} flag causes \LaTeX{} to write a \texttt{.fls} file
listing every file it reads.  The pcons helper converts this to a Ninja
depfile so that changes to any included \texttt{.tex}, \texttt{.sty},
or \texttt{.cls} file trigger a rebuild.


\section{Results}
\label{sec:results}

Cross-references\index{cross-reference} require multiple \LaTeX{} passes.
The table of contents\index{table of contents} also requires a second pass
to resolve page numbers.  As noted in Section~\ref{sec:intro}, pcons handles
all of this via latexmk.

\bibliographystyle{plain}
\bibliography{refs}

\printindex

\end{document}
